% GENERATED FILE. Edit canonical lessons, then run npm run generate:books.
% canonical-source-hash: fnv1a64:44c72f8ac0a4cbe1
% canonical-lessons: GU-C279-rit, GU-C279-kevi-rite-javu, GU-R279-asking-the-way

\chapter{Asking the way}
\label{ch:gu-asking-the-way}
\hlchaptermodality{279}

\begin{chapteropening}
\textbf{By the end of this chapter:} \emph{I can ask where a place is, how far it is and how to get there (sṭeśan kevī rīte javũ?), ask someone to show me the way, and follow and give simple directions.}

Complete GU-R279-asking-the-way: ask the way, and give it.
\end{chapteropening}

\section[\texorpdfstring{\gu{રીત} (rīt)}{રીત (rīt)}]{\gu{રીત} (rīt) — a way; a manner}

\label{lesson:GU-C279-rit}



\begin{quote}
Before the new one: say \textquotedblleft{}it was raining, so we stayed at home\textquotedblright{}. (\emph{varsād paḍto hato, eṭle ame ghare rahyā}.) Then \textquotedblleft{}by bus\textquotedblright{}. (\textbf{\gu{બસમાં}}.)
\end{quote}

\subsection*{You'll want to know: \gu{રીત}}
\textbf{\gu{રીત}} — \emph{rīt} — \textquotedblleft{}a way; a manner\textquotedblright{}: the way a thing is done. It is feminine: \textbf{\gu{તેની} \gu{રીત}}, \textquotedblleft{}his way\textquotedblright{}.

\begin{quote}
\textbf{\gu{તેની} \gu{રીત} \gu{સારી} \gu{છે}.}
\end{quote}

\begin{quote}
\emph{tenī rīt sārī che.} — \textquotedblleft{}His way of doing it is good.\textquotedblright{}
\end{quote}

\textbf{\gu{કેવી} \gu{રીતે}} (\emph{kevī rīte}), \textquotedblleft{}in what way\textquotedblright{}, is how Gujarati says \textquotedblleft{}how\textquotedblright{}: \textbf{\gu{કેવી}} asks \textquotedblleft{}what kind of\textquotedblright{}, agreeing with \textbf{\gu{રીત}}, and \textbf{\gu{રીતે}} takes the \textbf{-\gu{ે}} of \textbf{\gu{સવારે}}.

\begin{quote}
\textbf{\gu{તે} \gu{કેવી} \gu{રીતે} \gu{આવી}?}
\end{quote}

\begin{quote}
\emph{te kevī rīte āvī?} — \textquotedblleft{}How did she come?\textquotedblright{}
\end{quote}

The answer can be \textbf{\gu{બસમાં}}: by bus.

\begin{cousinweb}
\textbf{\gu{રીત}} (\emph{rīt}) goes back to Sanskrit \emph{rīti}, \textquotedblleft{}going, a course\textquotedblright{}, and later \textquotedblleft{}a way of doing things\textquotedblright{} (Monier-Williams). CDIAL's entry for \emph{rīti} (CDIAL 10751) lists a Gujarati \emph{rī}, \textquotedblleft{}method, manner\textquotedblright{}; \textbf{\gu{રીત}} keeps the Sanskrit \emph{-t}, so it is probably the learned form.
\end{cousinweb}

\subsection*{Your turn}
Say these aloud:

\begin{itemize}
\raggedright
  \item \emph{rīt}, then \emph{kevī rīte}
  \item \textquotedblleft{}How did she come?\textquotedblright{} — \emph{te kevī rīte āvī?}
  \item \textquotedblleft{}How did you come?\textquotedblright{}, with respect — \emph{tame kevī rīte āvyā?}
\end{itemize}

\begin{tcolorbox}[breakable,colback=teal!4,colframe=teal!35!black,title={Before you move on}]
What does \textbf{\gu{રીત}} mean? (\textquotedblleft{}A way; a manner\textquotedblright{}.) \textquotedblleft{}How?\textquotedblright{} (\textbf{\gu{કેવી} \gu{રીતે}?}) Why \textbf{\gu{કેવી}}, not \textbf{\gu{કેવો}}? (\textbf{\gu{રીત}} is feminine.)
\end{tcolorbox}
\section[\texorpdfstring{\gu{સ્ટેશન} \gu{કેવી} \gu{રીતે} \gu{જવું}? (sṭeśan kevī rīte …)}{સ્ટેશન કેવી રીતે જવું? (sṭeśan kevī rīte …)}]{\gu{સ્ટેશન} \gu{કેવી} \gu{રીતે} \gu{જવું}? (sṭeśan kevī rīte javũ?) — asking the way}

\label{lesson:GU-C279-kevi-rite-javu}



\begin{quote}
Say \textquotedblleft{}how?\textquotedblright{} (\emph{kevī rīte?}) Then \textquotedblleft{}where\textquotedblright{} and \textquotedblleft{}far\textquotedblright{}. (\textbf{\gu{ક્યાં}}, \textbf{\gu{દૂર}}.) And \textquotedblleft{}excuse me\textquotedblright{}. (\textbf{\gu{માફ} \gu{કરો}}.)
\end{quote}

\begin{grammarlens}[title={Grammar Lens: three questions}]
\begin{quote}
\textbf{\gu{માફ} \gu{કરો}, \gu{સ્ટેશન} \gu{કેવી} \gu{રીતે} \gu{જવું}?}
\end{quote}

\begin{quote}
\emph{māph karo, sṭeśan kevī rīte javũ?} — \textquotedblleft{}Excuse me, how do I get to the station?\textquotedblright{}
\end{quote}

The verb stays in its \textbf{-\gu{વું}} form: \textquotedblleft{}how to go to the station?\textquotedblright{} Three questions do most of the work: where it is, \textbf{\gu{સ્ટેશન} \gu{ક્યાં} \gu{છે}?}; how far, \textbf{\gu{કેટલું} \gu{દૂર} \gu{છે}?}; and how, \textbf{\gu{સ્ટેશન} \gu{કેવી} \gu{રીતે} \gu{જવું}?} One more asks to be shown: \textbf{\gu{મને} \gu{રસ્તો} \gu{બતાવશો}?} It is a polite future, \textquotedblleft{}will you show me the way?\textquotedblright{}: a request, not an order.
\end{grammarlens}

\subsection*{You'll want to know: the answer}
\begin{quote}
\textbf{\gu{સીધા} \gu{જાઓ}, \gu{પછી} \gu{પુલ} \gu{પાસે} \gu{ડાબી} \gu{બાજુ} \gu{વળો}.}
\end{quote}

\begin{quote}
\emph{sīdhā jāo, pachī pul pāse ḍābī bāju vaḷo.} — \textquotedblleft{}Go straight, then turn left at the bridge.\textquotedblright{}
\end{quote}

Directions are said with \textbf{-\gu{ઓ}}, the form of \textbf{\gu{માફ} \gu{કરો}}: \textbf{\gu{જાઓ}}, \textbf{\gu{વળો}}. \textbf{\gu{ડાબી}} and \textbf{\gu{જમણી}} follow \textbf{\gu{બાજુ}}, which is feminine. Other pieces: \textbf{\gu{બેંકની} \gu{સામે}} (\emph{beṅknī sāme}), \textquotedblleft{}opposite the bank\textquotedblright{}; \textbf{\gu{સિગ્નલ} \gu{પછી}} (\emph{signal pachī}), \textquotedblleft{}after the signal\textquotedblright{}; \textbf{\gu{દૂર} \gu{નથી}, \gu{દસ} \gu{મિનિટ} \gu{છે}} (\emph{dūr nathī, das miniṭ che}), \textquotedblleft{}it's not far, ten minutes\textquotedblright{}.

\subsection*{Your turn}
Say these aloud:

\begin{itemize}
\raggedright
  \item \textquotedblleft{}Will you show me the way?\textquotedblright{} — \emph{mane rasto batāvśo?}
  \item \textquotedblleft{}It's opposite the bank\textquotedblright{} — \emph{beṅknī sāme che}
\end{itemize}

\begin{tcolorbox}[breakable,colback=teal!4,colframe=teal!35!black,title={Before you move on}]
\textquotedblleft{}How far is it?\textquotedblright{} (\textbf{\gu{કેટલું} \gu{દૂર} \gu{છે}?}) \textquotedblleft{}How do I get to the station?\textquotedblright{} (\textbf{\gu{સ્ટેશન} \gu{કેવી} \gu{રીતે} \gu{જવું}?}) \textquotedblleft{}Turn left\textquotedblright{}? (\textbf{\gu{ડાબી} \gu{બાજુ} \gu{વળો}}.)
\end{tcolorbox}
\section[(review)]{Review — asking the way}

\label{lesson:GU-R279-asking-the-way}



\begin{quote}
Say \textquotedblleft{}you should go to the doctor\textquotedblright{}. (\emph{tamāre dāktar pāse javũ joīe}.) Then \textquotedblleft{}so\textquotedblright{}, the everyday way. (\emph{eṭle}.)
\end{quote}

\subsection*{Your turn}
Say these aloud:

\begin{itemize}
\raggedright
  \item Excuse me, where is the station? How far is it?
  \item How do I get to the station? Will you show me the way?
  \item Go straight, then turn left at the bridge.
  \item It's not far, ten minutes. It's opposite the bank.
  \item How did you come? By bus, because the train was late. Now turn it round, with eṭle.
\end{itemize}

\begin{tcolorbox}[breakable,colback=teal!4,colframe=teal!35!black,title={Before you move on}]
\textquotedblleft{}How?\textquotedblright{} (\textbf{\gu{કેવી} \gu{રીતે}?}) Which form gives a direction? (\textbf{-\gu{ઓ}}: \textbf{\gu{જાઓ}}, \textbf{\gu{વળો}}.) Why \textbf{\gu{ડાબી} \gu{બાજુ}}? (\textbf{\gu{બાજુ}} is feminine.)
\end{tcolorbox}
